% Level 2, Week 1 -- answer key
\answersweek{1}{Welcome Back, Python}{Answers for workbook chapter 1 and homework 1.}

\section{Workbook exercises}
\begin{answer}[Exercise 1.1 Hero Card]
\textbf{Step 1 (error):} \texttt{TypeError: can only concatenate str (not "int") to str}.\par
\textbf{Step 4 (types):} \texttt{hero\_name} str, \texttt{power\_points} int, \texttt{health} float, \texttt{is\_legend} bool.\par
\textbf{What to look for:} the level is converted with \texttt{int()} once, where it is typed in (not in every line);
every card line uses an f-string; the student explains why \texttt{health} has \texttt{.0} at the end (int times float gives a float).
Accept any labels and border style.
\end{answer}
\pyfile{../excercises/week01/ex-hero-card/solution.py}

\begin{answer}[Exercise 1.2 Mood Meter]
\textbf{Prediction table:} 9 + yes: Super charged.\quad 9 + no: Ready to play.\quad 2 + yes: Ready to play.\quad 3 + no: Recharge time.\par
\textbf{What to look for:} the \texttt{and} condition comes \emph{before} the \texttt{or} condition. If a student puts the
\texttt{elif} first, energy 9 with a snack wrongly says \enquote{Ready to play}: a great moment to ask
\enquote{which check does Python do first?} The tip uses \texttt{if not had\_snack:}.
\end{answer}
\pyfile{../excercises/week01/ex-mood-meter/solution.py}

\section{Homework}
\begin{answer}[Core: Type detective]
\texttt{42} int,\quad \texttt{"42"} str,\quad \texttt{4.2} float,\quad \texttt{False} bool,\quad \texttt{"False"} str,\quad \texttt{-8} int.
\end{answer}

\begin{answer}[Core: Dragon Door]
Ticks: key yes + level 5 opens; key yes + level 2 shut; key no + level 12 opens.
Also correct: \texttt{return (has\_key and level >= 3) or level >= 10} on one line.
\end{answer}
\pyfile{../excercises/week01/hw-dragon-door/solution.py}

\begin{answer}[Challenge: Predict, then check]
\texttt{False}, \texttt{True}, \texttt{False}, \texttt{77}, \texttt{14}.
\end{answer}

\begin{answer}[Creative: Design a hero]
Any made-up hero is fine. Check for one value of each type, for example
name \texttt{"Comet Kai"} (str), level \texttt{6} (int), speed \texttt{2.5} (float), can fly \texttt{True} (bool).
If a student writes a real name or school, gently remind them of the hero-name rule.
\end{answer}

\section{Placement quiz}
The placement quiz answers and marking guide are in the \textbf{Assessments} section at the end of this Answer Key.
